\documentclass[a4paper]{article}
\usepackage[a4paper,margin=1in]{geometry}
\usepackage{graphicx} % Required for inserting images
\usepackage{placeins} % Provides \FloatBarrier to keep floats in the right section
\usepackage{float} % Provides [H] to place floats exactly here
\usepackage{subcaption} % Side-by-side figures with subcaptions
\usepackage{indentfirst}
\usepackage{amsmath}
\usepackage{amssymb}
\usepackage{booktabs} 
\usepackage[backend=biber,style=numeric,sorting=none]{biblatex}
\addbibresource{references.bib}

% Reduce overfull lines (e.g., long arXiv/DOI strings) without changing margins.
\setlength{\emergencystretch}{2em}

\title{Evaluating LLM Agents in Black-Box Auctions}
\author{Nelson Nguyen, Chuan Xiao, Shuyuan Zheng}
\date{February 2026}

\begin{document}

\maketitle



\begin{abstract}
This study aims to evaluate the bidding performance of LLM agents in MenuNet auctions. Because MenuNet guarantees incentive compatibility, we measure the truthfulness and bidder regret of our agents across different auction conditions, comparing LLM agents with varying reasoning effort, models, and prompts. The LLM agents are prompted to place bids that maximize their internal utility, whether those bids are truthful or not. We find that all models perform well in auctions without history, whereas performance becomes more varied in auctions with history. Higher-reasoning models typically outperform lower-reasoning models, but the design of MenuNet introduces nuanced problems in the decision making of higher-reasoning models. Overall, this study demonstrates the ability of LLM agents as bidders in MenuNet auctions, while limitations point to future research on improving agent reasoning.
\end{abstract}



\section{Introduction}
Advances in technology have led to the development of neural networks in auction mechanisms due to their ability to learn under multiple constraints and handle high-dimensional auctions environments \cite{dutting2022optimal}. We see this in modern economics; for example, Alibaba has applied deep reinforcement learning to sponsored-search real-time bidding auctions \cite{zhao2018ssrtb}. More recently, large language models (LLMs), such as GPT-5, have advanced in reasoning and decision modeling, motivating their use as automated bidders in auctions. LLM agents provide a general approach to complex auctions where bidders must adapt to changing situations and adjust to multiple bidding objectives while quickly returning new bids.


Prior work has analyzed LLM performance in first-price, second-price, and RegretNet mechanisms; we extend upon these studies by examining LLM agents in MenuNet auctions \cite{anon2024blackbox}. While RegretNet can model multi-bidder settings, it typically provides only approximate incentive compatibility. In contrast, MenuNet guarantees incentive compatibility in a single-bidder setting by presenting a menu of allocation-payment options and allowing the buyer to select the preferred option.As a result, truthful reporting is optimal

We therefore study whether LLM agents in MenuNet auctions can both maximize their own utility while following the mechanism’s incentives by bidding truthfully. We evaluate multiple GPT models across reasoning levels and prompt variants, measuring truthfulness and bidder regret under three auction parameter settings in hopes of providing insight into the reliability of LLM agents as the overall applicability of both LLM agents as well as MenuNet auctions. To our knowledge, we are the first to evaluate LLM agents in MenuNet auctions.



\section{MenuNet Implementation}
MenuNet is a neural-network auction composed of two modules, Mechanism Network and Buyer Network. Given a bid report, it returns a single menu option consisting of allocation probabilities and the associated payment. Our implementation follows the architecture proposed in "Automated Mechanism Design via Neural Networks" \cite{shen2021amdnns}.

MechNet parameterizes a menu through a learnable allocation matrix $X \in [0,1]^{m \times k}$ and payment vector $p \in \mathbb{R}^{k}$, where $m$ is the number of items and $k$ is the number of menu options (including an exit option). We learn unconstrained parameters, $X_{\text{params}} \in \mathbb{R}^{m \times (k-1)}$ and $p_{\text{params}} \in \mathbb{R}^{k-1}$, then apply a sigmoid nonlinearity to obtain allocations (kept within [0,1]) and $\mathrm{softplus}$ to obtain payments (ensure non-negative). The exit option is added by appending a final zero column to $X$ and a final zero entry to $p$, so each column $X_{:,j}$ and payment $p_j$ define one menu option.

BuyerNet models option selection by evaluating utilities over a discretized valuation grid. Given upper bounds $\bar v$ for each item and $d=100$ discretization points per item, BuyerNet forms a Cartesian-product grid over $[0,\bar v_1] \times \cdots \times [0,\bar v_m]$ (we use $\bar v = [1,2]$). For each valuation profile $v$ on this grid and each menu option $j$, it computes
\[
U(v,j) = v^\top X_{:,j} - p_j
\]
Here, $v \in \mathbb{R}^m$ denotes the buyer's valuation vector over the $m$ items, $j \in \{1,\dots,k\}$ indexes a menu option, $X_{:,j} \in [0,1]^m$ is the allocation-probability vector for option $j$, $p_j \in \mathbb{R}_{\ge 0}$ is its payment, and $U(v,j)$ is the resulting quasi-linear utility ($U(v,j)$ = expected value - payment). These utilities are converted to a stochastic choice distribution via a temperature-scaled softmax: $S(v)=\mathrm{SoftMax}(U(v)/0.05)$ over $j$.


MenuNet computes expected revenue under a uniform distribution over the valuation grid as $\mathbb{E}_{v}\left[\sum_j S(v,j)\,p_j\right]$. We train for 5000 steps with Adam (learning rate 0.01) to minimize loss, the negative of revenue (equivalent to maximizing revenue), saving the best-performing parameters during training. At inference time, given a reported bid vector $v$, MenuNet returns the menu option $j^\star=\arg\max_j (U(v,j))$ and outputs $(X_{:,j^\star}, p_{j^\star})$.



\section{Auction Implementation}

We evaluate a single LLM bidder in two-item MenuNet auctions using valuation domains $[0,1] \times [0,2]$ and $k=3$ menu options. For each round, the bidder's private valuations are sampled independently and uniformly from these intervals unless we enable a same-bidder setting, in which case the bidder's valuation vector is held fixed across rounds. Each auction round consists of the LLM agents submitting a bid vector, the trained MenuNet returning a menu option, and the bidder receiving its utility.

We test bidder behavior under three experimental factors: (i) \emph{history} vs.\ \emph{no history}, (ii) \emph{rule} vs.\ \emph{no rule}, and (iii) \emph{same bidder} vs.\ \emph{different bidders}. Under the history condition, the prompt includes prior rounds (valuation domains, bids, and returned menu options) as well as a listing of all available menu options for each valuation domain encountered; without history, each round is presented independently. Under the rule condition, the prompt additionally provides a high-level description of the mechanism's internal computation (allocation matrix, payment vector, and utility-based option selection) without explicitly naming MenuNet; without rule, the prompt describes only the bidding task and utility. The same-bidder condition keeps valuations fixed across rounds, while the different-bidder condition resamples valuations each round.

We compare models and reasoning effort settings supported by the OpenAI Responses API, varying the reasoning effort parameter (minimal, low, medium) and using two prompt variants: one that requests only bids and one that requests a brief rationale together with bids. When history is enabled, rounds are queried sequentially because the prompt grows with the round history; we truncate history depending on the reasoning setting to reduce runtime. Unless otherwise noted, we run 30-round auctions per condition and repeat conditions across multiple runs to obtain at least 30 rounds of data per setting.

To quantify behavior, we measure (i) \emph{truthfulness}, defined as bidding exactly the valuation vector (to the same decimal precision), and (ii) \emph{regret}, defined as the difference between the profit induced by the submitted bid and the profit induced by truthful bidding whenever the two reports select different menu options.

\begin{table}[ht]
\centering
\begin{tabular}{lccc}
\toprule
Allocation & Item 1 & Item 2 & Payment\\
\midrule
Menu Option 1        & 0.9995  & 0.9996 & 1.3357 \\
Menu Option 2        & 0.9993  & 0.0171  & 0.6760 \\
Menu Option 3 (exit) & 0  & 0  & 0 \\
\bottomrule
\end{tabular}
\caption{ Menu Options Returned by MenuNet Trained on Domains [0,1] * [0,2] }
\end{table}



\section{Results}
We report \emph{truthfulness} as the number of rounds in which the LLM agent's bids exactly matches its item valuations and \emph{regret} as the cumulative profit difference between the submitted bid and truthful bidding summed over all rounds. Negative regret indicates that deviating from truthful bidding reduced profit.


\subsection{GPT-5-mini: Prompt 1 vs Prompt 2}
\subsubsection{Prompt 1}
\begin{figure}[H]
\centering
\includegraphics[width=0.98\linewidth]{prompt1_avg_by_preset.png}
\caption{Average truthfulness and regret by auction preset for Prompt 1 across reasoning effort settings.}
\label{fig:prompt1_avg_by_preset}
\end{figure}

Prompt 1 requests only a bid tuple each round (no explicit rationale). Across all reasoning-effort settings, performance without history is near-perfect: the agent bids truthfully in essentially all rounds and incurs negligible regret. When history is enabled, we find the agent has more trouble maintaining optimal bidding. Truthfulness decreases and becomes sensitive to auction conditions and reasoning effort (Figure~\ref{fig:prompt1_avg_by_preset}).

While we may expect higher reasoning models to perform better, we surprisingly find that minimal reasoning performs best in history auctions with \emph{different bidders}, averaging near perfect truthfulness and zero regret. However, when switching to \emph{same bidder}, the inverse occurs and medium reasoning performs best, also averaging near perfect truthfulness and zero regret (Table~\ref{tab:prompt1_history_by_preset}). Notably, median regret is zero across Prompt 1 settings, indicating that most deviations from truthful bidding do not change the selected menu option.

\begin{table}[H]
\centering
\begin{tabular}{lrrrr}
\toprule
Reasoning & \multicolumn{2}{c}{\emph{different bidders}} & \multicolumn{2}{c}{\emph{same bidder}} \\
\cmidrule(lr){2-3}\cmidrule(lr){4-5}
 & Truth & Regret & Truth & Regret \\
\midrule
\texttt{minimal} & 29.91 & -0.008 & 27.18 & -0.161 \\
\texttt{low}     & 26.07 & -0.023 & 28.48 & -0.004 \\
\texttt{medium}  & 22.03 & -0.132 & 29.63 & -0.014 \\
\bottomrule
\end{tabular}
\caption{Prompt 1, history-only: average truthfulness (out of 30) and cumulative regret (over 30 rounds), pooled over \emph{rule} and \emph{no rule}.}
\label{tab:prompt1_history_by_preset}
\end{table}



\FloatBarrier
\subsubsection{Prompt 2}
\begin{figure}[H]
\centering
\includegraphics[width=0.98\linewidth]{prompt2_avg_by_preset.png}
\caption{Average truthfulness and regret by auction preset for Prompt 2 across reasoning effort settings.}
\label{fig:prompt2_avg_by_preset}
\end{figure}

Prompt 2 requests both a brief rationale and the bid tuple each round. Like Prompt 1, performance without history is also near-perfect. With history enabled, Prompt 2 leads to substantially worse regret for minimal reasoning effort, marginally worse regret in low reasoning effort, but significant improvements for medium reasoning effort agent (Figure~\ref{fig:prompt_compare_truth_regret}). Note, 

In medium reasoning models, although regret is significantly reduced in \emph{different bidders} auctions, we find that truthfulness has also dropped. This is paradoxical to the nature of MenuNet, where truthful bids are the optimal bids. However, we find that models with reasoning will purposely bid dishonestly in order to guarantee MenuNet returns the option they consider most profitable. As a result, truthfulness drops but regret improves.


\begin{table}[H]
\centering
\begin{tabular}{lrrrr}
\toprule
Reasoning & \multicolumn{2}{c}{\emph{different bidders}} & \multicolumn{2}{c}{\emph{same bidder}} \\
\cmidrule(lr){2-3}\cmidrule(lr){4-5}
 & Truth & Regret & Truth & Regret \\
\midrule
\texttt{minimal} & 29.80 & -0.032 & 20.59 & -2.437 \\
\texttt{low}     & 18.43 & -0.057 & 25.70 & -0.068 \\
\texttt{medium}  & 18.65 & -0.024 & 29.75 &  0.000 \\
\bottomrule
\end{tabular}
\caption{Prompt 2, history-only: average truthfulness (out of 30) and cumulative regret (over 30 rounds), pooled over \emph{rule} and \emph{no rule}.}
\label{tab:prompt2_history_by_preset}
\end{table}



\FloatBarrier
\subsubsection{Overall Results}
Across all reasoning effort, removing history yields near-perfect performance for both prompts: truthfulness is approximately $30/30$ and regret is essentially zero across rule and bidder settings. In addition, LLM Agents peform slightly better in all auctions with the rule parameter on when compared to rule off, but \emph{rule} does not cause the the same variance in LLM performance as \emph{history} or \emph{same bidder}.

With history enabled, we observe large truthfulness drops  when valuations change across rounds (\emph{different bidders}), where Prompt 2 low/medium reasoning averages roughly $18$--$19$ truthful rounds and Prompt 1 medium reasoning averages roughly $21$--$23$ truthful rounds out of $30$ (Figures~\ref{fig:prompt1_avg_by_preset} and~\ref{fig:prompt2_avg_by_preset}). We also observe a consistent gap between truthfulness and regret: even when bids are not exactly truthful, the selected menu option is often unchanged, yielding near-zero regret. This suggests that MenuNet can be robust to small deviations in reported bids, but it also highlights that truthfulness alone may overstate economically meaningful deviations. We see this when comparing medium reasoning effort between Prompt 1 and Prompt 2 in \emph{different bidder} auctions with \emph{history}, in which although truthfulness dropped significantly, regret also reduced significantly (Figures~\ref{fig:prompt_compare_truth} and~\ref{fig:prompt_compare_regret}). Conversely, the heavy-tailed regret behavior observed for Prompt 2 with minimal reasoning shows that some prompt and reasoning configurations can trigger misbidding that changes menu selection and substantially reduces utility (Figure~\ref{fig:history_truncation_ablation}).

\begin{figure}[ht]
\centering
\begin{subfigure}[t]{0.49\linewidth}
  \centering
  \includegraphics[width=\linewidth]{compare_prompt_compare_truth.png}
  \caption{Truthfulness.}
  \label{fig:prompt_compare_truth}
\end{subfigure}
\hfill
\begin{subfigure}[t]{0.49\linewidth}
  \centering
  \includegraphics[width=\linewidth]{compare_prompt_compare_regret.png}
  \caption{Regret.}
  \label{fig:prompt_compare_regret}
\end{subfigure}
\caption{History-only auction comparison of Prompt 1 vs Prompt 2 across auction presets, grouped by reasoning effort.}
\label{fig:prompt_compare_truth_regret}
\end{figure}

\begin{figure}[ht]
\centering
\includegraphics[width=0.98\linewidth]{compare_history_truncation_ablation.png}
\caption{Minimal-reasoning ablation controlling for history truncation: mean truth rate and mean regret for Prompt 1 vs Prompt 2.}
\label{fig:history_truncation_ablation}
\end{figure}



\subsection{Other Models}

\begin{figure}[H]
\centering
\begin{subfigure}[t]{0.49\linewidth}
  \centering
  \includegraphics[width=\linewidth]{v3_avg_by_preset.png}
  \caption{Average performance by preset.}
  \label{fig:v3_avg_by_preset}
\end{subfigure}
\hfill
\begin{subfigure}[t]{0.49\linewidth}
  \centering
  \includegraphics[width=\linewidth]{v3_figure_1.png}
  \caption{Reasoning-level comparison.}
  \label{fig:v3_reasoning_compare}
\end{subfigure}
\caption{Additional comparisons including GPT-4o and GPT-5.2 across auction presets. }
\label{fig:v3_model_comparisons}
\end{figure}

While other GPT models were also tested, data collection was not as rigorous as in GPT-5-mini; auctions for these models were only run using Prompt 2. Figure~\ref{fig:v3_model_comparisons} indicates that auction performance is highly model-dependent in history settings. GPT-4o performs poorly under history, with very low truthfulness and large negative regret, suggesting that it frequently selects suboptimal menu options. GPT-5.2 improves substantially over GPT-4o, but it still trails GPT-5-mini variants in the most challenging presets, particularly when valuations change across rounds (\emph{different bidders}). By contrast, GPT-5-mini remains comparatively robust across presets, typically achieving high truthfulness and near-zero regret. Similar to medium reasoning effort in GPT-5-mini, we find deviations from exact truthfulness in Prompt 2 often reflect strategic misreports that preserve the selected menu option rather than economically meaningful mistakes. In fact, GPT-5.2 appears to misreport more often, where in history auctions, GPT-5.2 consistently has a lower average truthfulness but maintains similar regret to GPT-5-mini using Prompt 1.



\section{Conclusion}

Across all no-history settings, performance is near-perfect regardless of model, prompt, or reasoning effort: regret is essentially zero and bids are almost always exactly truthful. This indicates that LLM agents can act as reliable bidders when each round is presented independently. However, many real-world auctions are not memory-less; bidders often remember or observe past outcomes, so the more practically relevant regime is the history setting, where behavior becomes substantially more variable.

In the more practical situation where history is stored, LLM agents tend to struggle. For most cases, we observe higher reasoning models perform more truthfully and have less regret in MenuNet auctions. However, only for \emph{same bidder} auctions using Prompt 1, the inverse occurs, where the higher reasoning model performs significantly worse than the minimal reasoning model. The stark contrast in performance between for GPT-5-mini between Prompt 1 and Prompt 2 reveals LLM agent's sensitvity to prompt design, where even minor changes will lead to major differences in results.

In the more practical setting where round history is available, LLM bidding behavior becomes markedly less stable. In most conditions, higher reasoning effort improves performance, yielding higher truthfulness and lower regret. A notable exception occurs in \emph{same bidder} history auctions only under Prompt 1, where higher reasoning performs worse than minimal reasoning; one possible explanation is increased reasoning leads to "overthinking", resulting in avoidable deviations from the dominant truthful strategy. More broadly, the gap between GPT-5-mini’s outcomes under Prompt 1 and Prompt 2 highlights strong sensitivity to prompt design, where relatively small changes in instruction framing can produce large shifts in measured truthfulness and regret.

Generally, we observe the trend that agents that bid more truthfully tend to have lower regret and thus perform better. This aligns will with MenuNet being incentive compatible. However, we see nuances in this as the medium reasoning effort models have less \emph{truthfulness} yet also lower \emph{regret}. LLM agent's at times bid over their internal valuations if they deem Menu Option 1 as the best option. While this violates both individual rationality and incentive compatability, it demonstrates LLM agent's complex thinking as well as ability to recognize the mechanisms of the auctions and create "loopholes" in order to maximize their utility. Whether this behavior is desired is up to implementation, and can likely be prevented by adjusting the prompt.



\subsection{Discussion and Future Work}
Our experiments consider a simple single-bidder environment (two items, three menu options, fixed valuation domains). A natural extension is to scale the mechanism (more items and menu options) and expand valuation domains to test whether performance degrades as the decision space grows. We also plan to study settings with \emph{different items} across rounds, where the valuation domain changes over time together with the bidder's sampled valuations, which better reflects practical online auction environments.

Beyond scaling, future work should identify the mechanisms behind the observed model and prompt dependence in history settings. In particular, GPT-5.2 underperforms GPT-5-mini in several history presets, which may reflect sensitivity to prompt format or adherence to the intended input/output constraints. Finally, we can investigate cases where higher-reasoning agents intentionally overbid relative to their valuations to force a preferred menu option, and whether such behavior is and desirable from a mechanism-design perspective.



\newpage
\printbibliography[title={references}]
\end{document}
